% Lapisan solusi O001 -- Bab 7 (Operator Kompak)
% 1 solusi asli terpisah; bukan tulisan John M. Erdman.
% Provenans model: OpenAI Codex gpt-5.6-sol, Ultra, atas arahan pengguna.
% Lisensi komponen solusi: CC BY-SA 4.0.
% Tidak menyiratkan dukungan John M. Erdman atau Portland State University.

\markboth{SOLUSI O001 UNTUK BAB 7: OPERATOR KOMPAK}{SOLUSI O001 UNTUK BAB 7: OPERATOR KOMPAK}
\section*{Solusi O001 untuk Bab 7: Operator Kompak}
\addcontentsline{toc}{section}{Solusi O001 untuk Bab 7: Operator Kompak}
\begin{center}
\small
Komponen solusi asli terpisah, CC BY-SA 4.0. Ditulis dengan bantuan
OpenAI Codex gpt-5.6-sol, Ultra, atas arahan pengguna. Pernyataan latihan
berasal dari edisi terjemahan karya John M. Erdman; solusi ini bukan tulisan
beliau dan tidak menyiratkan dukungan beliau atau Portland State University.
\end{center}

% O001-SOLUTION-ID: O001-FAOA-2015-CH07-EX-001-SOLUTION
% SOURCE-EXERCISE-ID: FAOA-2015-CH07-NODE-0045
% STATEMENT-TARGET-SHA256: d34a802dca82cf381a646022b8b8cec4c40d52cfbb09ae67d5a13d39e1d48396
\begin{o001solution}
  {O001-FAOA-2015-CH07-EX-001-SOLUTION}
  {FAOA-2015-CH07-NODE-0045}
  {d34a802dca82cf381a646022b8b8cec4c40d52cfbb09ae67d5a13d39e1d48396}
\begin{o001statement}
Misalkan $(S, \fml A,\mu)$ suatu ruang ukur $\sigma$-hingga dan $\phi \in
L_\infty(\mu)$. Tentukan dekomposisi polar dari
 \index{operator perkalian!dekomposisi polar dari}%
 \index{operator!perkalian!dekomposisi polar dari}%
 \index{dekomposisi polar!dari operator perkalian}%
operator perkalian $M_\phi$ (lihat contoh~\ref{C063527}).
\end{o001statement}
\begin{o001answer}
Pada $L_2(S,\mu)$, misalkan $E=\{s:|\phi(s)|>0\}$ hingga himpunan nol dan
definisikan
\[
 u(s)=\begin{cases}\phi(s)/|\phi(s)|,&s\in E,\\0,&s\notin E.\end{cases}
\]
Maka
\[
  |M_\phi|=M_{|\phi|},\qquad M_\phi=M_uM_{|\phi|},
\]
dan $M_u$ adalah isometri parsial dengan proyeksi awal sekaligus proyeksi
akhir $M_{\mathbf1_E}$.
\end{o001answer}
\begin{o001proof}
Dari Contoh~\ref{C063527}, ${M_\phi}^*=M_{\overline\phi}$. Oleh karena itu
\[
 {M_\phi}^*M_\phi=M_{|\phi|^2}
 \quad\hbox{dan}\quad
 |M_\phi|=({M_\phi}^*M_\phi)^{1/2}=M_{|\phi|},
\]
karena $M_{|\phi|}$ positif dan kuadratnya $M_{|\phi|^2}$.

Fungsi $u$ terukur, terbatas secara esensial, dan
$|u|^2=\mathbf1_E$ hampir di mana-mana. Maka
\[
  M_u^*M_u=M_{|u|^2}=M_{\mathbf1_E},
\]
sebuah proyeksi ortogonal; jadi $M_u$ isometri parsial. Perhitungan titik
demi titik memberi
\[
 M_uM_{|\phi|}=M_{u|\phi|}=M_\phi.
\]
Kernel $M_\phi$ terdiri atas kelas fungsi yang nol hampir di mana-mana pada
$E$, sehingga komplemen ortogonalnya adalah $L_2(E)$, yakni jangkauan
$M_{\mathbf1_E}$. Ini ruang awal $M_u$.

Untuk memeriksa ruang akhir yang disyaratkan dekomposisi polar, ambil
$g\in L_2(E)$ dan tetapkan
$g_n=g\mathbf1_{\{|\phi|\ge1/n\}}$. Maka $g_n\to g$ dalam $L_2$ melalui
konvergensi terdominasi, sedangkan
\[
 g_n=M_\phi\left(\frac{g_n}{\phi}\right)
 \quad\hbox{dan}\quad
 \frac{g_n}{\phi}\in L_2.
\]
Jadi $\overline{\ran M_\phi}=L_2(E)$. Karena $M_uM_u^*=M_{\mathbf1_E}$,
ruang akhir $M_u$ juga $L_2(E)$. Semua syarat dekomposisi polar, termasuk
kasus $\phi=0$ hampir di mana-mana, telah terpenuhi.
\end{o001proof}
\end{o001solution}
